\section{Discussion}
\label{sec:discussion}

\subsection{Relation to Vol.\ 1}

Vol.~1 established that time and space emerge from the causal log $C_{ij}$
through the accumulation of closure density $\sigma$.
The present volume shows that the internal structure of matter ---
particle types, symmetries, charges, and masses ---
can be read from the same closure geometry.

The key bridge is the closure energy functional~\eqref{eq:closure_energy}.
In Vol.~1, this functional governed the emergence of metric structure.
In the present volume, its local minima are identified with particle types,
and its symmetry classes with gauge groups.
The two volumes thus share a common geometric foundation.

\subsection{Scope and Limitations}

The interpretations advanced in this volume are structural and qualitative.
Several points of caution are warranted.

First, no quantitative predictions are derived.
The mass relation $m \propto 1/\Delta^2$ and the charge formula $q_i = \phi_i - 1/3$
are structural analogies, not derived equalities.
Precise numerical correspondence with observed values
would require a more developed formalism.

Second, the identification of closure regimes with particle classes is proposed
as a consistent mapping, not as a proof of equivalence.
The claim is that non-degenerate triangular closure \emph{admits}
the structure of baryonic states, not that it \emph{is} a baryon.

Third, the derivation of $\mathrm{SU}(3)$ from non-degenerate closure
identifies it as a \emph{candidate} consistent with closure invariance,
not as a unique derivation.
A rigorous uniqueness argument would require additional constraints.

\subsection{Path toward Vol.\ 3}

Vol.~3 addresses the large-scale consequences of closure geometry:
the emergence of gravity and spacetime curvature
from closure density, vorticity, and metric distortion.
The central question is whether the Einstein equations
can be read as a fixed-point condition of the closure network.

The present volume provides the particle-physics half of this picture.
The two halves --- internal structure (Vol.~2) and large-scale geometry (Vol.~3) ---
are expected to be unified within the same closure framework.
